\section{Read a fixed header before interpreting chunks}
The first fourteen bytes are a header. The magic is the four ASCII bytes for
\texttt{qoif}. Width and height are unsigned 32-bit integers in big-endian order.
The channels byte is three or four. The colorspace byte is zero or one. Zero denotes
sRGB color components with linear alpha; one denotes linear components. These two
fields describe the image; they do not select a different chunk grammar. In particular,
do not use channels equal to three as a reason to misclassify an RGBA chunk.

\tagpdfsetup{table/header-rows={1}}
\begin{center}
  \begin{tabular}{rll}
    Offset & Bytes                & Meaning                  \\
    \hline
    0      & \texttt{71 6f 69 66} & Magic                    \\
    4      & \texttt{00 00 00 08} & Width eight              \\
    8      & \texttt{00 00 00 01} & Height one               \\
    12     & \texttt{04}          & Four-channel description \\
    13     & \texttt{00}          & sRGB with linear alpha   \\
  \end{tabular}
\end{center}

A byte load is independent of the pointer's alignment. Read the four dimension bytes
as unsigned values and assemble the numerical integer with shifts of twenty-four,
sixteen and eight bits. Do not load a native little-endian word and treat its value
as the file's dimension. Validate the dimension before multiplying it by the other
one. The local teaching domain is width and height from one through sixty-four, so
the product is at most 4096 pixels. Those limits are exercise limits, not limits of
the QOI format.

After the header, the next input offset is fourteen and the next output index is zero.
The dimensions determine how many pixels must be reconstructed. Neither a zero byte
nor the first occurrence of the eventual end-marker pattern determines that count.
